% =====================================================================
%  04-documentos.tex - Repositorio documental y firma individual
% =====================================================================

\section{Repositorio documental y firma individual}\label{sec:impl-documentos}

Cada documento pertenece a un expediente y conserva su contenido cifrado, sus versiones inmutables y una clasificación con historia propia. Esta sección recorre las operaciones documentales, la admisión de nuevas cargas y de soportes PDF y DOCX, las consultas de metadatos, el versionado, la clasificación, la entrega del contenido original con su buzón de incidentes y la firma individual con custodia de la clave por su titular.

\subsection{Operaciones documentales por expediente}\label{sec:impl-http-documentos}

La carga recibe hasta 16~MiB binarios y un nombre seguro en \texttt{X-Document-Name}. El actor deriva de \texttt{Authorization: Bearer}; \texttt{X-Actor} no forma parte del contrato. El servicio crea UUID y versión uno, calcula SHA-256 y prepara el paquete cifrado \texttt{DVLT1}. El AAD conserva el par identidad documental--versión; la asociación con el expediente es una restricción inmutable de PostgreSQL, no un campo agregado retroactivamente al contexto criptográfico.

\begin{table}[H]
  \centering
  \caption{Contrato documental por expediente. Todas las rutas de las seis primeras filas usan el prefijo \texttt{/api/v1/cases/<case>/documents}.}
  \label{tab:impl-http-rutas}
  \begin{tabularx}{\linewidth}{@{}L{5.2cm} X@{}}
    \toprule
    \textbf{Método y sufijo} & \textbf{Responsabilidad} \\
    \midrule
    \texttt{POST} al prefijo & Carga cifrada y asociada al expediente; \texttt{201}. \\
    \texttt{GET} al prefijo & Lista metadatos con paginación y filtros por nombre y estado de sellado. \\
    \texttt{GET /<document>} & Devuelve metadatos persistidos del documento en el expediente indicado. \\
    \texttt{POST /<document>/seal} & Firma con la clave del servidor y sello local; requiere una sola versi\'on disponible. \\
    \texttt{POST /<document>/verify} & Verificación del documento con una sola versión disponible. \\
    \texttt{GET /<document>/evidence} & ZIP del documento con una sola versión; cabecera \texttt{X-Document-Digest}. \\
    \midrule
    \texttt{GET /api/v1/audit/verify} & Owner comprueba la cadena PostgreSQL completa. \\
    \texttt{GET /healthz} & Diagnóstico público del proceso. \\
    \bottomrule
  \end{tabularx}
\end{table}

Owner accede a cualquier expediente. Litigante y Paralegal necesitan pertenencia vigente para consultar metadatos, clasificar, cargar, a\~nadir versiones, verificar o exportar. El sellado con clave del servidor se conserva para Owner y Litigante; la firma individual de la \cref{sec:impl-firma-individual} corresponde exclusivamente al Litigante con su propio v\'inculo. Cliente conserva denegadas todas las acciones documentales incluso cuando est\'a asignado. El rol se comprueba antes de revelar el recurso. Un documento fuera del expediente indicado responde \texttt{404 document\_not\_found}, incluso para Owner o un usuario perteneciente a ambos expedientes. Las antiguas rutas globales bajo \texttt{/api/v1/documents} responden \texttt{404} y no tienen alias.

Carga, detalle y sellado conservan los campos \texttt{case\_id}, \texttt{id}, \texttt{version}, \texttt{name}, \texttt{digest} y \texttt{sealed}. Listado, detalle actual, carga y anexado agregan \texttt{current\_metadata} con la clasificaci\'on vigente; detalle hist\'orico, historial de contenido y sellado conservan una proyecci\'on exclusivamente documental. Cada versi\'on conserva su propio contenido cifrado y estado de sellado. La ruta \texttt{seal} descifra y comprueba el digest, produce y verifica la firma con la clave configurada del servidor, solicita el token y lo verifica bajo la cadena TSA capturada. Una segunda solicitud por esa ruta sobre un documento sellado produce \texttt{409 document\_already\_sealed}. La firma individual presenta pruebas producidas en el equipo del titular y dispone de su propio recibo recuperable. Exportar comprueba el contenido descifrado y empaqueta la evidencia capturada sin actualizarla; la consulta de verificaci\'on proporciona por separado el reporte criptogr\'afico completo.

\subsection{Admisión general de nuevas cargas documentales}\label{sec:impl-admision-documental}

La carga binaria (\cref{tab:impl-http-rutas}), la carga con clasificación
inicial (\cref{sec:impl-http-clasificacion}) y el anexado de versiones
(\cref{sec:impl-http-versiones}) comparten una política obligatoria de admisión
antes del hash, cifrado y registro. El límite es de 16~MiB por archivo. Se admiten ocho familias: PDF, DOCX,
TXT, JPEG, PNG, MP3, WAV y MP4. El nombre, la extensión, el MIME declarado y los
metadatos no seleccionan el formato ni permiten omitir su validación. El
contenido admitido conserva exactamente sus bytes originales; la inspección no
lo transforma en otra versión ni modifica documentos, firmas o sellos históricos.

PDF y DOCX reutilizan los analizadores de la política de soportes procesales
(\cref{sec:impl-formatos}). TXT exige UTF-8 válido y
admite BOM inicial, tabulación y saltos de línea, pero rechaza otros controles.
JPEG y PNG deben ser imágenes estáticas completas. MP3 requiere Layer~III; WAV
admite PCM de 8, 16, 24 o 32 bits, hasta ocho canales y 192~kHz. MP4 admite
contenedores regulares con H.264 o AAC, hasta ocho pistas, sin fragmentación,
cifrado ni referencias externas de datos. Cada cuadro de imagen o video queda
limitado a 16\,777\,216 píxeles. Para multimedia se comprueban la estructura, el
inventario permitido y la decodificación completa de todas las pistas a salida
nula; reconocer una cabecera o leer una muestra no basta para aceptar el archivo.

La aplicación comprueba identidad y acceso antes de inspeccionar. La inspección
no mantiene abierta una transacción de persistencia. Antes de confirmar se
reautentica al principal completo; la transacción documental conserva la
comprobación de pertenencia, cierre administrativo, versión esperada y auditoría.
Un rechazo no cifra ni persiste una nueva versión. Qadra distingue formato no
admitido, formato inválido, límite de validación e indisponibilidad del validador,
y conserva archivo y nombre junto con la clasificación de la carga o la versión
esperada del anexado hasta una acción explícita, sin reenvío automático
(\cref{sec:impl-qadra-conciliacion}). Cliente mantiene denegado el acceso
documental (\cref{sec:impl-http-documentos}).

La entrada acotada se copia a un descriptor \texttt{memfd} sellado contra
escritura, crecimiento y reducción. Los procesos secuenciales reciben ese
mismo contenido mediante descriptor. La inspección dispone de cinco segundos
de CPU, 256~MiB de espacio de direcciones y diez segundos transcurridos. La
sonda y la decodificación disponen de dos y ocho segundos de CPU, respectivamente,
y 512~MiB de espacio de direcciones cada una. Las cuotas suman quince segundos
de CPU y comparten un vencimiento monótono de veinte segundos transcurridos;
no se reinician por pista ni se transfieren entre procesos. Se acota también la
salida del protocolo y se fijan en cero las cotas de escritura de archivos y
volcados de memoria.

El proceso privado instala los límites y cierra descriptores heredados ajenos
antes de ejecutar el decodificador. El supervisor conserva la identidad del
líder hasta terminar su grupo y recolectar el hijo directo. La implementación
nativa usa FFmpeg~9.0.2 con componentes restringidos y red deshabilitada en su
compilación; verifica las ocho familias al arrancar. Desconectar el cliente
HTTP no cancela inmediatamente una validación ya iniciada: ésta conserva el
supervisor y su presupuesto. Estas medidas no constituyen un aislamiento
completo del sistema operativo ni un antivirus, y las páginas del descriptor
pueden llegar a swap. Un archivo válido que exceda recursos se rechaza por
límite, sin presentarlo como corrupción.

Esta admisión general es independiente de la política de soportes procesales
PDF/DOCX de la \cref{sec:impl-formatos}: admitir audio, video o imágenes como
documentos no los habilita donde un acto exige uno de esos soportes. Tampoco
acredita autenticidad jurídica ni amplía permisos. El contrato preciso reside en
\rutarepo{docs/document-upload-admission-api.md} y su justificación en
\rutarepo{docs/adr/0058-bounded-general-document-admission.md}. Su verificación
y aceptación se presentan en la \cref{sec:pruebas-admision-documental}.


\subsection{Admisión acotada de PDF y DOCX}\label{sec:impl-formatos}

{\raggedright La política \texttt{pdf\_docx\_v1}, descrita en \rutarepo{docs/adr/0024-isolated-document-format-admission.md}, comprueba el contenido de nuevos soportes procesales. No se aplica retrospectivamente ni declara resuelta la validación de todas las cargas generales. Cada lote admite hasta dos referencias distintas, 16~MiB por documento y 32~MiB en total. El adaptador limita vault y evidencia antes de transferirlos; la evidencia JSON admite 1~MiB antes de deserializarse. La selección definitiva repite las cotas en su propia instantánea.\par}

La aplicación presta los bytes ya autenticados a un proceso separado del mismo ejecutable. Un protocolo binario acota entradas y respuesta; el texto del documento no pasa por argumentos ni mensajes de error. El worker Linux instala 5 segundos de CPU, 256~MiB de espacio de direcciones, tamaño de archivo de salida cero y core dumps deshabilitados antes de leer documentos. El supervisor impone 10 segundos desde el lanzamiento, usa tuberías no bloqueantes y recoge al proceso al fallar. Al concluir, libera buffers y parsers, vacía la respuesta y termina sin handlers \texttt{atexit}; el hijo no escribe perfiles LLVM que infringirían el límite de archivos. Las pruebas de biblioteca permanecen instrumentadas. El presupuesto pertenece al lote completo; no se renueva por archivo ni incluye criptografía o consultas del padre.

PDF usa la API C de qpdf~12.4.1, cargada dentro del worker desde una instalación con artefacto y SHA-256 fijados. Deshabilita recuperación, rechaza advertencias, errores y cifrado, incluso con contraseña de apertura vacía. Comprueba catálogo, referencias, padres y contadores del árbol de páginas, sin ciclos, con máximo de 4096 páginas y profundidad acotada. Esta política no promete un contador de expansión PDF, renderizado, validación completa del estándar ni detección de malware.

DOCX admite un perfil ZIP32 con Stored o Deflate y hasta 1024 entradas físicas. Comprueba cabeceras, directorio, nombres inequívocos, descriptores, CRC32, tamaño real y final exacto del flujo comprimido. Cada entrada se expande una vez; los DOCX del lote comparten 64~MiB de expansión y un millón de eventos XML, con profundidad 128. Exige tipos de contenido, relaciones internas, parte principal Word y un cuerpo; admite familias Transitional y Strict coherentes. Rechaza DTD, entidades personalizadas, macros, OLE, ActiveX, paquetes embebidos, plantillas adjuntas y \texttt{altChunk}. Solo admite relaciones externas de hipervínculo inertes HTTP, HTTPS o mailto, que no se siguen.

{\raggedright La admisión no extrae archivos ni resuelve enlaces; tampoco es una validación XSD completa de OPC/OOXML. Puede rechazar documentos válidos fuera de este perfil, como XML UTF-16 o ZIP64. Se diferencian formato rechazado, tamaño excesivo y presupuesto agotado. Un fallo del parser no autoriza la transición ni reescribe evidencia. El servidor comprueba la instalación nativa antes de abrir su puerto; \rutarepo{scripts/setup-document-formats.sh} prepara la dependencia y \rutarepo{docs/document-format-operations.md} documenta operación, límites y reproducción.
\par}
\subsection{Consultas autorizadas de metadatos}\label{sec:impl-http-consultas}

El listado recibe \texttt{limit}, \texttt{offset}, \texttt{name}, \texttt{sealed}, \texttt{document\_type}, \texttt{classification} y \texttt{tag}. El tamaño de página predeterminado es 50 y admite de 1 a 100; el desplazamiento es un entero no negativo. El filtro de nombre admite hasta 200 caracteres, recorta espacios exteriores y rechaza caracteres de control. PostgreSQL busca una subcadena literal sin distinguir mayúsculas mediante \texttt{strpos(lower(name), lower(...))}; los caracteres \texttt{\%} y \texttt{\_} conservan su significado literal. El filtro de sellado distingue documentos pendientes y documentos con evidencia capturada. La búsqueda alcanza el nombre almacenado, no el contenido cifrado.

El adaptador restringe primero al expediente autorizado y selecciona la versión de mayor número de cada UUID y su última revisión de clasificación; después aplica los filtros y ordena por UUID antes de paginar. Tipo, categoría y una etiqueta individual se comparan por igualdad exacta sensible a mayúsculas y se combinan mediante conjunción. Una etiqueta retirada en una revisión posterior deja de coincidir; comas, acentos y caracteres de comodín se conservan literalmente. Una versión histórica sellada no hace aparecer como sellada a la versión actual pendiente. Obtiene como máximo una fila adicional para calcular \texttt{has\_more}; la respuesta contiene \texttt{documents} y ese indicador, sin un conteo global. El orden es estable para un conjunto sin modificaciones; varias peticiones no comparten una instantánea, por lo que altas concurrentes pueden modificar la composición de las páginas. Un expediente visible sin coincidencias devuelve una página vacía.

La proyección SQL obtiene identidad, versión, nombre, digest, presencia de evidencia y valores de clasificación; no carga ni descifra el vault y no deserializa firmas, certificados o tokens. El detalle actual reúne la versión máxima de la asociación expediente/documento y la clasificación vigente en una misma transacción; el detalle histórico conserva únicamente el contenido del número solicitado. Por ello, obtener metadatos no constituye una verificación criptográfica de integridad: el usuario solicita esa comprobación mediante la operación específica de verificación, que presenta sus cuatro componentes.

El permiso \texttt{ReadDocument} habilita las consultas para Owner, Litigante y Paralegal; los dos últimos requieren pertenencia vigente. Cliente permanece denegado incluso si pertenece al expediente. Cada consulta autentica la sesión y la transacción PostgreSQL revalida cuenta activa, rol y pertenencia antes de leer. El listado registra \texttt{document.listed}, incluso cuando no hay coincidencias; el detalle registra \texttt{document.read}. Los metadatos se entregan solamente después de confirmar su evento de auditoría. Un fallo de inserción o commit devuelve error sin liberar el resultado de la consulta.

Los contratos residen en \rutarepo{crates/application/src/documents/query.rs} y \rutarepo{crates/web/src/routes/queries.rs}; la proyección persistida se implementa en \rutarepo{crates/infrastructure/src/document_postgres/query.rs}. Estas operaciones conservan los documentos cifrados y la evidencia existente sin modificar su versión, firma o sello.

\subsection{Versiones inmutables y selección de evidencia}\label{sec:impl-http-versiones}

El tipo \texttt{Document\allowbreak Version\allowbreak Ref} identifica una instant\'anea mediante UUID y n\'umero positivo de versi\'on. \texttt{Document\allowbreak Processor} prepara cada contenido nuevo con una DEK aleatoria propia y el AAD de ese par, utilizando el formato \texttt{DVLT1} existente. Dos versiones con bytes iguales comparten digest, pero solicitan cifrados independientes. La firma separada del PDF y el token RFC~3161 cubren el digest del contenido; por s\'i mismos no firman el UUID, el n\'umero de versi\'on, el expediente ni el nombre. En la firma individual, el consentimiento adicional liga expediente, UUID y versi\'on con la cuenta y las pruebas presentadas. En el sellado con clave del servidor, el actor de auditor\'ia deriva de la sesi\'on y no se equipara al titular de esa clave.

La raíz \texttt{document\_series} fija UUID, expediente y \texttt{first\_available\_version}. Las filas de \texttt{documents} usan clave primaria compuesta \texttt{(id, version)}; las claves foráneas impiden asociaciones incompatibles y una restricción diferida exige que la primera versión exista al confirmar. La versión actual se obtiene como \texttt{MAX(version)}, sin un puntero mutable independiente. Si se importa únicamente la versión 7, esa es la primera disponible y la siguiente es 8; el sistema no inventa las versiones 1 a 6. Nombre, digest, vault y contexto no se modifican. La evidencia pasa de ausente a capturada una sola vez por versión.

Anexar requiere \texttt{expected\_version}. La aplicación autoriza y compara la cabeza antes de cifrar; tras preparar, autentica nuevamente. El almacén vuelve a comprobar permisos, expediente y cabeza dentro de la transacción auditada. Dos solicitudes que esperan la misma versión producen un solo sucesor y un conflicto para la restante. Una cabeza distinta devuelve \texttt{409 document\_version\_conflict}; el agotamiento del contador de 32 bits devuelve \texttt{409 document\_version\_exhausted}. El evento \texttt{document.version\_added} y la fila se confirman juntos; fallar cualquiera de ellos impide el resultado exitoso.

{\emergencystretch=3em Las rutas de la \cref{tab:impl-http-versiones} utilizan el prefijo común
\rutarepo{/api/v1/cases/<case>/documents/<document>/versions}.\par}

\begin{table}[H]
  \centering
  \caption{Contrato de consulta y operación sobre versiones documentales.}
  \label{tab:impl-http-versiones}
  \begin{tabularx}{\linewidth}{@{}L{5.4cm} X@{}}
    \toprule
    \textbf{Método y sufijo} & \textbf{Responsabilidad} \\
    \midrule
    \texttt{POST}\newline\texttt{?expected\_version=<n>} & Añade contenido binario con \texttt{X-Document-Name}; devuelve \texttt{201} y la nueva versión pendiente. \\
    \texttt{GET ?limit=<n>\&}\newline\texttt{before\_version=<v>} & Historial descendente de metadatos con cursor exclusivo y página de 1 a 100 elementos. \\
    \texttt{GET /<version>} & Detalle de la versión indicada en el expediente. \\
    \texttt{POST /<version>/seal} & Firma con la clave del servidor y sello de esa versi\'on; conserva las restantes. \\
    \texttt{POST /<version>/verify} & Reporte criptográfico con UUID, expediente y versión evaluada. \\
    \texttt{GET /<version>/evidence} & ZIP histórico; cabeceras \texttt{X-Document-Id}, \texttt{X-Document-Version} y \texttt{X-Document-Digest}. \\
    \bottomrule
  \end{tabularx}
\end{table}

{\emergencystretch=3em El historial devuelve la lista \texttt{versions} y el indicador \texttt{has\_more}. Incluye el cursor \texttt{next\_before\_version} y la primera versión disponible, \texttt{first\_available\_version}. La página siguiente solicita versiones menores que la última entregada; anexar una nueva cabeza no desplaza las páginas ya recorridas. Para verla se refresca la primera página. Historial y detalle consultan exclusivamente metadatos y confirman sus eventos antes de responder. Un número inexistente y una asociación ajena producen el mismo error de recurso ausente. Cliente permanece denegado; Owner, Litigante y Paralegal conservan las condiciones de acceso al expediente, y Paralegal no puede sellar.\par}

Las acciones con versi\'on expl\'icita resuelven una sola instant\'anea y conservan ese par durante preparaci\'on y confirmaci\'on. Las rutas de sellado, verificaci\'on y exportaci\'on sin n\'umero funcionan \'unicamente si hay una versi\'on disponible al resolver; si hay varias, responden \texttt{409 document\_version\_required}. Si una operaci\'on ya resolvi\'o la \'unica versi\'on y despu\'es se a\~nade otra, confirma la instant\'anea original sin redirigirse a la nueva cabeza. Sellar una versi\'on anterior y a\~nadir un sucesor son operaciones compatibles. La verificaci\'on general conserva la pol\'itica y el instante de evaluaci\'on de su contrato; la consulta individual agrega el contexto hist\'orico capturado y separa el estado actual del certificado. Conservar bytes no garantiza que una TSA contin\'ue siendo aceptada indefinidamente.

{\emergencystretch=3em Los contratos y la resolución se encuentran en \rutarepo{crates/application/src/documents/version.rs}, \rutarepo{crates/application/src/documents/case_content.rs} y \rutarepo{crates/web/src/routes/versions.rs}. El fundamento y las limitaciones se documentan en \rutarepo{docs/adr/0019-immutable-document-versions.md}.\par}

\subsection{Clasificación documental con historia independiente}\label{sec:impl-http-clasificacion}

El valor \texttt{Document\allowbreak Metadata}, definido en \rutarepo{crates/domain/src/document_metadata.rs}, representa tipo documental, categoría manual y etiquetas. Los dos campos opcionales admiten hasta 80 caracteres Unicode cada uno; los espacios exteriores se recortan y un resultado vacío se interpreta como ausencia. Se rechazan caracteres de control sobre la entrada original, antes del recorte. La solicitud admite hasta 20 etiquetas antes de deduplicar, cada una de 1 a 40 caracteres. Las etiquetas se ordenan por bytes UTF-8 y se deduplican por igualdad exacta. Se conservan mayúsculas, acentos, formas Unicode, comas y espacios interiores; no hay inferencia automática ni cambio de permisos a partir de estos valores.

La clasificación pertenece a la raíz documental, no a una versión de archivo. La revisión cero significa que nunca se registró una decisión y devuelve campos ausentes y etiquetas vacías, sin autor ni fecha inventados. Las filas persistidas empiezan en uno. Cada reemplazo aceptado añade la revisión siguiente, incluso si limpia los campos o repite los valores; no sobrescribe historia. El contador llega hasta $2^{32}-1$ y su agotamiento es un conflicto explícito. El tipo \texttt{Metadata\allowbreak Revision} se mantiene separado de \texttt{Document\allowbreak Version}.

El canon de valores empieza con \texttt{DMETA1}. Cada campo opcional usa un indicador de ausencia o presencia; los valores presentes se codifican como longitud de cuatro bytes en orden de red y bytes UTF-8. Sigue la cantidad de etiquetas y cada etiqueta con su longitud y contenido. El máximo es 3\,940 bytes. El puerto SHA-256 calcula su huella, y PostgreSQL contrasta la misma representación con su función nativa, sin extensión adicional. Esta huella identifica valores organizativos: no cubre UUID, autor, revisión ni instante y no equivale a una firma de la decisión. La identidad documental y la revisión forman parte del recurso auditado; autor e instante se conservan por separado.

La tabla \texttt{document\_metadata\_revisions} conserva una clave compuesta por UUID y revisión, campos normalizados, digest, instante UTC y UUID y correo histórico del actor. Ambos datos de identidad proceden del mismo usuario recargado y bloqueado dentro de la transacción, por lo que consultar historia no reconstruye el correo mediante la cuenta vigente. Los valores están en claro y son visibles únicamente bajo el permiso de lectura documental del expediente. Decodificar una fila comprueba canon, huella e instante; un dato inconsistente provoca error, sin normalización silenciosa.

{\emergencystretch=3em Las rutas de la \cref{tab:impl-http-clasificacion} usan el prefijo \rutarepo{/api/v1/cases/<case>/documents}. La operación de reemplazo recibe \texttt{expected\_\allowbreak metadata\_\allowbreak revision}, tipo, categoría y el vector de etiquetas, y rechaza campos JSON desconocidos.\par}

\begin{table}[H]
  \centering
  \caption{Contrato HTTP de clasificación documental.}
  \label{tab:impl-http-clasificacion}
  \begin{tabularx}{\linewidth}{@{}L{5.2cm} X@{}}
    \toprule
    \textbf{Método y sufijo} & \textbf{Responsabilidad} \\
    \midrule
    \texttt{POST /with-metadata} & Carga multipart de archivo y clasificación inicial; confirma ambos y devuelve \texttt{201}. \\
    \texttt{GET /<document>/}\newline\texttt{metadata} & Estado actual con UUID, expediente, revisión y valores; no fabrica procedencia para la revisión cero. \\
    \texttt{PUT /<document>/}\newline\texttt{metadata} & Reemplazo completo condicionado a la revisión esperada; devuelve la revisión confirmada. \\
    \texttt{GET /<document>/}\newline\texttt{metadata/history} & Historia descendente con \texttt{limit} de 1 a 100 y cursor exclusivo \texttt{before\_revision}. \\
    \bottomrule
  \end{tabularx}
\end{table}

El permiso \texttt{Classify\allowbreak Document} corresponde a Owner, Litigante y Paralegal dentro de su ámbito documental; Cliente permanece denegado. La transacción toma el bloqueo común de auditoría y revalida cuenta, rol, pertenencia y asociación documental antes de comparar la revisión esperada. Una cabeza diferente devuelve \texttt{document\_\allowbreak metadata\_\allowbreak conflict}; el agotamiento, \texttt{document\_\allowbreak metadata\_\allowbreak revision\_\allowbreak exhausted}, ambos con estado 409. El cambio y \texttt{document.metadata\_changed} se confirman juntos. Lectura e historia registran sus propios eventos antes de responder, aun cuando el resultado sea revisión cero o una página vacía.

La carga clasificada recibe exactamente una parte \texttt{file} y una parte JSON \texttt{metadata}, en cualquier orden. El nombre autorizado sigue en \texttt{X-Document-Name}; el nombre del formulario multipart no sustituye ese contrato. El archivo admite 16~MiB, JSON 8~KiB y la petición completa 16~MiB más 32~KiB. El adaptador consume primero el cuerpo completo bajo ese límite, incluidos bytes posteriores al cierre multipart, y rechaza un fallo de transporte antes de preparar contenido. Después valida partes completas, únicas y conocidas, límites individuales y valores. La aplicación exige permisos de carga y clasificación antes y después de cifrar y vuelve a autenticar a la misma identidad. Raíz, versión inicial, revisión uno y los dos eventos de éxito se confirman en una única transacción. La carga binaria existente conserva el estado sin clasificación; la carga multipart explícita crea revisión uno aunque sus campos opcionales estén vacíos.

\texttt{Document\allowbreak Overview} reúne contenido y clasificación actual y se devuelve desde la propia transacción de carga o anexado, sin un GET posterior que pudiera fallar después del éxito. Reclasificar no compara la versión del archivo; anexar, sellar, verificar o exportar una versión no compara la revisión organizativa. La clasificación no entra en \texttt{Document\allowbreak Record}, AAD, \texttt{DVLT1}, firma, sello, certificados ni ZIP de evidencia. Las decisiones y límites se fundamentan en \rutarepo{docs/adr/0020-audited-document-classification.md}.

\subsection{Contenido original e incidentes de integridad}\label{sec:impl-contenido-documental}

Además de verificar y exportar evidencia (\cref{tab:impl-http-versiones}), el repositorio entrega el contenido original de una versión exacta. La consulta de contenido identifica expediente, UUID documental y versión positiva. Recupera los bytes originales tanto si están pendientes de sello como si ya contienen evidencia capturada; no sustituye la verificación de firma, certificado o sello de tiempo. Una versión nueva no redirige la petición de una anterior. Owner accede a todos los expedientes y Litigante o Paralegal requieren pertenencia vigente. Cliente queda denegado; el cierre administrativo conserva esta lectura autorizada.

El servicio carga una captura cifrada autorizada, valida el sobre, descifra mediante AES-GCM con el AAD de identidad y versión y compara SHA-256. Después reautentica al principal completo. La transacción final comprueba cuenta, rol, pertenencia y coincidencia del contenido inmutable, y confirma el evento \texttt{document.content\_authorized}. Sólo entonces el transporte recibe los bytes. El evento acredita autorización confirmada, no recepción por el navegador ni escritura en disco. Redis y PostgreSQL no forman una transacción distribuida; cerrar sesión después de liberar los bytes no puede recuperarlos.

La ruta admite hasta 16~MiB de contenido. PostgreSQL comprueba el tamaño cifrado antes de materializarlo y repite la cota en la selección; el procesador aplica su propio límite. El cupo de descargas acompaña al buffer ceroizable hasta liberar la última referencia de bytes, incluso si el manejador HTTP ya terminó; un receptor lento no libera prematuramente esa capacidad. Un histórico mayor devuelve error de tamaño sin eliminarse ni generar un incidente de integridad. Consultar tampoco somete una carga histórica a la admisión de la \cref{sec:impl-admision-documental}. La respuesta satisfactoria usa adjunto binario, nombre saneado, longitud exacta, prohibición de caché y cabeceras de expediente, identidad, versión y digest. No se entregan rangos parciales: la validación comprende el archivo completo.

Un fallo de validación produce una transacción independiente que conserva el incidente y \texttt{document.content\_rejected}; no crea un acceso satisfactorio. Registra identidades exactas, solicitante histórico, clase del fallo, instantes y huellas, sin texto claro, vault, tokens ni claves. El actor técnico identifica la observación y no acusa al solicitante. El UUID de observación permite repetir exactamente el registro sin duplicarlo; reutilizarlo con otros datos produce conflicto. Intentos distintos mantienen incidentes distintos.

Las clases separan sobre inválido, autenticación criptográfica fallida, digest discordante y cambio de captura. No determinan la causa: una KEK incorrecta también puede causar fallo de autenticación. Un error de base de datos o auditoría sigue siendo técnico; si no se confirma el incidente, no se afirma que Owner haya sido notificado. El rechazo conserva cero bytes de archivo en la respuesta.

La tabla \texttt{document\_integrity\_incidents} es inmutable y forma parte del catálogo, inventario y respaldo. Sólo Owner activo consulta su buzón, con autorización y auditoría actuales; el orden paginado es UUID ascendente, no cronológico. Qadra consulta existencia al entrar, permite actualizar y abre la versión exacta tras revalidar acceso. La apertura no descarga automáticamente contenido rechazado. El aviso no acredita correo, reparación ni atención humana. Como en las demás vistas (\cref{sec:impl-qadra-conciliacion}), la descarga descarta resultados tardíos al cambiar expediente, versión o sesión. El contrato se conserva en \rutarepo{docs/document-content-api.md} y la decisión en el ADR~0042 de \rutarepo{docs/adr/}.


\subsection{Firma individual con custodia externa}
\label{sec:impl-firma-individual}

La firma individual distingue la cuenta autenticada, el certificado p\'ublico y
la autorizaci\'on de una versi\'on documental. Solamente el rol Litigante puede
registrar su propio v\'inculo y presentar una firma personal; para firmar necesita
adem\'as pertenencia vigente al expediente y que este se encuentre activo.
Owner y Paralegal conservan sus permisos de consulta, verificaci\'on y
exportaci\'on, sujetos a las reglas de acceso por expediente, sin adquirir por
ello permiso para firmar como otro usuario. Cliente mantiene denegado el
acceso documental. La posesi\'on de un certificado no sustituye la sesi\'on
con MFA ni concede pertenencia.

El titular conserva su clave privada fuera del navegador y del servidor.
Despu\'es de generar una solicitud CSR y obtener el certificado de firma de la
CA interna, importa solamente material p\'ublico. El servicio
\texttt{Signing\allowbreak Binding\allowbreak Service} prepara una declaraci\'on
de 142 bytes con prefijo \texttt{DOCBIND1}, identidad de despliegue y cuenta,
contadores de la cuenta, identificador del v\'inculo, huellas de ra\'iz y
certificado y revisi\'on de confianza. El titular descarga esos bytes, los
firma externamente y presenta la firma separada. El registro comprueba la
prueba de posesi\'on, el perfil del certificado, su cadena y CRL, y confirma
el v\'inculo junto con su evento de auditor\'ia. El sujeto o correo del
certificado no se utiliza para inferir la cuenta. Retirar un v\'inculo impide
su uso posterior y conserva su registro original; no revoca por s\'i mismo
el certificado en la CA.

Para un documento, el cliente descarga la versi\'on exacta seleccionada y
contrasta sus identificadores y digest con los metadatos recibidos. El
servidor descifra y valida el contenido almacenado y exige que sea PDF.
La firma separada del PDF se produce externamente con RSA-3072 y SHA-256.
A partir de ella, \texttt{Document\allowbreak Signing\allowbreak Service}
prepara un consentimiento de 258 bytes con prefijo \texttt{DOCSIGN1}.
Adem\'as del contexto de cuenta, certificado y confianza, incorpora el
identificador de operaci\'on, expediente, UUID y n\'umero de versi\'on,
digest del contenido y digest de la firma del PDF. La segunda firma externa
autoriza esos bytes can\'onicos. Por tanto, dos versiones con contenido
id\'entico no comparten consentimiento, y sustituir la firma del PDF
invalida su correspondencia con la autorizaci\'on.

El env\'io final contiene el consentimiento y las dos firmas, sin clave,
contrase\~na de clave ni contenedor PKCS12. El servicio reconstruye el contexto
vigente y exige igualdad exacta con el consentimiento presentado. Verifica
ambas pruebas antes de solicitar el token RFC~3161 y usa
\texttt{Document\allowbreak Processor::seal\_external} para incorporar el
material p\'ublico del firmante. Una respuesta TSA incorrecta o ausente no
confirma una firma documental. La plataforma vuelve a autenticar la sesi\'on
despu\'es del trabajo criptogr\'afico; la transacci\'on revalida cuenta,
pertenencia, estado del expediente, v\'inculo, confianza y ventana temporal
despu\'es de adquirir los bloqueos de confirmaci\'on.

La aceptaci\'on conserva conjuntamente la evidencia, el recibo y el evento
de auditor\'ia. Su instante procede del reloj del servidor dentro de la
confirmaci\'on. El token TSA cubre el digest del PDF; no incluye el
consentimiento ni acredita el instante de aceptaci\'on personal. La
auditor\'ia y el recibo registran esa aceptaci\'on dentro de la frontera de
confianza de la plataforma. La TSA local conserva su alcance de demostraci\'on
t\'ecnica y no se presenta como constancia de un PSC autorizado.

El identificador de operaci\'on permite recuperar una respuesta perdida.
La consulta autorizada del recibo devuelve la aceptaci\'on original, y repetir
el mismo consentimiento con ambas firmas id\'enticas devuelve ese mismo
recibo. Si dos env\'ios concurrentes preparan tokens distintos, prevalece la
primera confirmaci\'on y el segundo recupera su evidencia sin reemplazarla.
Cambiar cualquiera de las pruebas bajo la misma operaci\'on produce conflicto.
La recuperaci\'on conserva la autorizaci\'on actual de lectura, pero no exige
que el certificado o la CRL originales sigan vigentes ni vuelve a pedir un
sello. Un cambio de cuenta, confianza o v\'inculo que afecte una nueva
operaci\'on requiere preparar y firmar un consentimiento nuevo.

El adaptador HTTP expone el registro bajo
\rutarepo{/api/v1/identity/signing-certificates} y la firma bajo
\begin{flushleft}
  \small
  \texttt{/api/v1/cases/<case>/documents/<document>/}\\
  \texttt{versions/<version>/individual-signing}.
\end{flushleft}
Las operaciones documentales separan preparaci\'on, env\'io, consulta del
recibo y verificaci\'on. Los cuerpos tienen l\'imite de 32~KiB, rechazan
campos desconocidos o repetidos y admiten material binario mediante Base64
can\'onico. Los identificadores del consentimiento deben coincidir con la
ruta. La verificaci\'on devuelve por separado la evidencia hist\'orica, el
estado actual del certificado y el retiro del v\'inculo. Los contratos
completos se conservan en \rutarepo{docs/http-individual-signing.md}.

